\documentclass[10pt]{article}

% graphicx package, useful for including eps, pdf, and raster graphics
\usepackage{graphicx}
\DeclareGraphicsExtensions{.png,.png,.jpg,.pdf}

% basic packages
\usepackage{fancyhdr}
\usepackage{color}
\usepackage{parskip}
\usepackage{float}
\usepackage{microtype}
\usepackage{url}
\usepackage{hyperref}
\usepackage{booktabs}
\usepackage{makecell}
\usepackage{multirow}
\usepackage{pbox}
\usepackage{enumitem}

% page layout
\usepackage[top=1in, bottom=1in, left=0.75in, right=0.75in, includefoot, heightrounded]{geometry}

% Useful for seeing the layout as you tweak it.
%\usepackage{showframe}

\setlength{\headheight}{1.1in}
\pagestyle{fancy}
\fancyhf{}

% fonts
\usepackage{fontspec}

\usepackage{polyglossia}
\setmainlanguage{amharic}
\setotherlanguage{english}
\newfontfamily\ethiopicfont{AbyssinicaSIL-R.ttf}


\setmainfont{DroidSans}[
  Extension   = .ttf,
  UprightFont = *,
  BoldFont    = *-Bold ]
\setmonofont{DroidSansMono}[
  Extension   = .ttf,
  UprightFont = * ]

\hypersetup{colorlinks,urlcolor=blue}
\urlstyle{same}


\lhead{\includegraphics[height=1in,clip]{logo}}
\rhead{\pbox[b][1in][c]{\textwidth}{
    \small
    \textbf{ምርመራው የተደረገው በ፡-}\\
    Northwest Genomics Center\\
    University of Washington School of Medicine\\
    Genome Sciences Box 355065\\
    3720 15th Ave NE\\
    Seattle, WA 98195-5065}}

\newcommand{\PageLine}{\rule{\textwidth}{0.25mm}}

\begin{document}

\begin{description}[font=\normalfont,align=left,labelwidth=8em]
\item [\textbf{\VAR{pat_name|e}}]
\item [\textbf{የትውልድ ቀን፡-}] \VAR{birth_date|e}
\item [\textbf{ናሙና፡-}] ከመካከለኛ አፍንጫ ፣ በራስ የተሰበሰበ/የተወሰደ
\item [\textbf{የስብስብ ባርኮድ፡}] \VAR{qrcode|e}
\item [\textbf{የስብስብ ቀን፡}] \VAR{collect_ts|e}
\item [\textbf{የሪፖርት ቀን፡}]
  %- if status_code == "never-tested"
  ተመርምሬ አላውቅም
  %- else
  \VAR{result_ts|e}
  %- endif
\end{description}

\PageLine

\begin{center}
\Large
\textbf{የምርምር ጥናት ውጤት - ይህንን ውጤት ከጤና እንክብካቤ አቅራቢዎ ጋር እንዲጋሩ ይበረታታሉ}
\end{center}

\bigskip

\large \underline{\textbf{ስለዚህ ምርመራ}}

\begin{itemize}
\item

  COVID-19 ማለት SARS-CoV-2 በሚባለው ቫይረስ የሚከሰት በሽታ ነው። COVID-19 እንዳለብዎት ለመወሰን
  ይህ የምርምር ጥናት ምርመራ ከአፍንጫዎ በተወሰደው ናሙና ውስጥ ቫይረሱን ይፈልጋል/ይመረምራል። ለዚህ ዓይነቱ
  ምርመራ የቴክኒክ ቃል ``ተቃራኒ-ግልባጭ የፖሊመሬዝ ሰንሰለት ምላሽ፣'' ወይም በአጭሩ ``RT-PCR'' ነው።

\end{itemize}

\bigskip

\large \underline{\textbf{አስፈላጊ ማስታወሻዎች}}

የምርምር ጥናቱን የምርመራ ውጤትዎን ከማየትዎ በፊት፣ የሚከተሉትን ማየት አለብዎት፡-

\begin{itemize}
\item

  የሰበሰቡት \textbf{የናሙናው ዓይነት} እና \textbf{የላብራቶሪው ምርመራ የምርምር ጥናት} አካል እና ይህ
  ሪፖርት የእርስዎን \textbf{የምርምር ጥናት ምርመራ ውጤት} ይገልጻል።

\item

  ይህንን ናሙና እራስዎ ወስደዋል እና በራስዎ (ወይም ለሌላ ሰው) ሰይመውታል እና ወደ ላብራቶሪው በመላኪያ
  አገልግሎት መልሰውታል። እንደ ነርስ ወይም ሀኪም ባለ ፈቃድ ባለው የጤና እንክብካቤ ባለሙያ አልተሰበሰበም፣
  አልተሰየመም ወይም አልተዘዋወረም። ይህ የምርመራ ውጤትዎ ላይ ተጽዕኖ ሊያሳድር ይችላል፣ እናም የጤና እንክብካቤ
  ባለሙያ እንደገና ሊመረምርዎት ይፈልግ ይሆናል።

\item

  በዚህ \textbf{የምርምር ጥናት} ውስጥ መሳተፍ የህክምና እንክብካቤዎን ቦታ አይወስድም። ጽኑ የህመም ስሜት
  ከተሰማዎት ወይም ስለምርመራ ውጤትዎ ሃሳብ ከያዞት፣ ሀኪም ወይም የጤና እንክብካቤ አቅራቢ ማግኘት አለብዎት።

\end{itemize}

\bigskip
\PageLine

\large \underline{\textbf{የ SARS-CoV-2 (COVID-19) የምርምር ጥናት ምርመራ ውጤት}}

%- if status_code == "negative"
አልተገኘም (አሉታዊ)\\
%- elif status_code == "positive"
ተገኝቷል (አዎንታዊ)\\
%- elif status_code == "inconclusive"
በቂ ያልሆነ\\
%- elif status_code == "never-tested"
ለመመርመር አልተቻለም\\
%- else
\VAR{0/0} %# poor man's raise
%- endif

\PageLine
\bigskip

\large \underline{\textbf{ውጤትዎን መረዳት}}

ይህ ሠንጠረዥ የእርስዎ ሙከራ ሊያሳየው የሚችለውን ውጤት ሁሉ ይሰጣል እንዲሁም እነዚህ ውጤቶች ምን ማለት
እንደሆኑ ያብራራል፡፡ ይህንን \textbf{የምርምር ጥናት ምርመራ ውጤት} ማተም ወይም በኢሜል መላክ ወይም
ለሀኪምዎ ወይም ለጤና እንክብካቤ አቅራቢዎ ማጋራት ይችላሉ።

የፈተናዎ ውጤት ምንም ይሁን ምን ቤትዎ ውስጥ መቆየት እና ህመም የሚሰማዎት ከሆነ ከሌሎች ሰዎች ጋር ከመሆን
ይቆጠቡ፡፡ \textbf{የበሽታ ምልክቶችዎ ከቀጠሉ ፣ በጣም ህመም ከተሰማዎት ወይም ስለ ህመምዎ ከተጨነቁ ወይም
ስለምርምር ጥናቱ የምርመራ ውጤትዎ ጥያቄዎች ካሉዎት ሀኪም ወይም የጤና እንክብካቤ አቅራቢን ማነጋገር አለብዎት
፡፡}

\begin{itemize}


\item

  \textbf{If your test says ``Not Detected (Negative)''}, it means our lab did not find the
  virus that causes COVID-19 in the sample that you sent to us. No test is
  perfect, and there is a small chance that you could have the virus even though
  your test is negative.

\item

  \textbf{If your test says ``Detected (Positive)''}, it means our
  lab found the virus that causes COVID-19 in the sample that you sent us. You
  should stay home and avoid being around other people, especially if you are
  feeling sick.

  Your result will be shared with the Washington State and local health
  departments. You do \textbf{not} need to report these results yourself. The
  health department may call you to ask questions about your illness and people
  you may have been around while you were sick.

\item

  \textbf{If your test says ``Inconclusive''}, it means our lab is not certain
  of your test result. No test is perfect, and an inconclusive result happens
  when we find a small amount of SARS-CoV-2 in your sample, but not enough for
  us to say that you have COVID-19.

\item

  \textbf{If your test says ``Unable to be tested''}, it means there was either
  not enough material on your sample for us to test, or there was an error when
  we ran the test. Since we are not able to test your sample, we could not find
  whether or not you have the virus that causes COVID-19.

\end{itemize}

\bigskip

\large \underline{\textbf{መርጃዎች}}

ስለ COVID-19 የበለጠ መረጃ ከፈለጉ እባክዎን ከህዝብ ጤና ባልደረባዎችዎ እነዚህን መርጃዎች ይመልከቱ፡

\begin{itemize}
\item
  Public Health \textemdash Seattle \& King County (የማህበረሰብ ጤና - ሲያትል እና ኪንግ
  ካውንቲ)\\
  \url{https://www.kingcounty.gov/depts/health/covid-19/languages/amharic.aspx}

\item

  Washington State Department of Health (የዋሺንግተን ስቴት የጤና ክፍል)፡-\\
  \url{http://www.doh.wa.gov/Emergencies/Coronavirus}

\item
  U.S. Centers for Disease Control and Prevention (CDC፣ የአሜሪካ የበሽታ ቁጥጥር እና
  መከላከል ማዕከላት)፡-\\
  \url{https://wwwn.cdc.gov/pubs/other-languages?Sort=Lang%3A%3Aasc&Language=Amharic}

\end{itemize}

በስልክ ከአንድ ሰው ጋር መነጋገር ከፈለጉ እባክዎን ከሚከተሉት አንዱን ይደውሉ፡-

\begin{itemize}

\item

  ከ COVID-19 ጋር የተዛመዱ የህክምና ጥያቄዎች፡-ከ 8AM እስከ 7PM መካከል ባለው ሰዓት የ King County
  Call Center (ኪንግ ካውንቲ ጥሪ ማእከልን) በ \textbf{206-477-3977} ያግኙ

\item

  \textbf{ስለ COVID-19 ህክምና-ያልሆኑ ጥያቄዎች፣ ማክበርን እና ከንግድ ጋር የተዛመዱ ጉዳዮችን ጨምሮ፡-} ከ
  King County COVID-19 Business and Community Information Line (ኪንግ ካውንቲ
  COVID-19 ንግድ እና ማህበረሰብ መረጃ መስመር) ጋር በ \textbf{206-296-1608፣} ከሰኞ - ዓርብ ከ8:30
  AM እስከ 4:30 PM ባለው ጊዜ ውስጥ ይገናኙ

\item

  \textbf{በዋሽንግተን ግዛት ውስጥ ስለ COVID-19 አጠቃላይ ጥያቄዎች፡-} የ Washington State Novel
  Coronavirus Call Center (ዋሽንግተን ግዛት ኖቬል ኮሮና ቫይረስ የጥሪ ማእከል)ን ከ 6 AM እና 10 PM
  ባሉት ሰዓታተ መካከል በ \textbf{1-800-525-0127} ያግኙ

\end{itemize}

\bigskip

\large \underline{\textbf{ዘዴዎች እና ገደቦች}}

\textbf{ዘዴዎች}\\
ይህ ሙከራ የተደረገው በተቃራኒ ግልባጩ PCR ን በመጠቀም SARs-CoV-2 ን ለመለየት ነው፡፡ ይህ ሙከራ የቫይረሱ
ክፍሎች ORf1B እና S ን ማጉላት የሚጠቀም ሲሆን እንዲሁም የሰው RnaseP ማጉያ መቆጣጠሪያን ፈልጎ ያገኛል።
ይህ ምርመራ የተሰራው እና የአፈፃፀም ባህሪው በ University of Washington Northwest Genomics
Center የተወሰነ ነው። ይህ ላቦራቶሪ ከፍተኛ ውስብስብነት ያላቸውን የክሊኒካል ላብራቶሪ ምርመራ ለማከናወን ብቁ
በመሆኑ በ  WA Department of Health የጤና ክፍል የህክምና ምርመራ ጣቢያ ፈቃድ ተሰጥቶታል።

\textbf{ገደቦች}

\begin{itemize}

\item

  እነዚህ ውጤቶች እንደ የምርምር ጥናት አካል ሆነው ቀርበዋል፡፡

\item

  ይህንን ናሙና እራስዎ ወስደዋል እና በራስዎ (ወይም ለሌላ ሰው) ሰይመውታል እና ወደ ላብራቶሪው በመላኪያ
  አገልግሎት መልሰውታል። እንደ ነርስ ወይም ሀኪም ባለ ፈቃድ ባለው የጤና እንክብካቤ ባለሙያ አልተሰበሰበም፣
  አልተሰየመም ወይም አልተዘዋወረም። ይህ የምርመራ ውጤትዎ ላይ ተጽዕኖ ሊያሳድር ይችላል፣ እናም የጤና እንክብካቤ
  ባለሙያ እንደገና ሊመረምርዎት ይፈልግ ይሆናል።

\item

  ይህ ሪፖርት በምርምር ጥናት ላይ የተመሠረተ ስለሆነ የእርስዎ ውጤቶች በ SARS-CoV-2 ኢንፌክሽን የሚከሰተውን
  COVID-19 በሽታን፣ ለይቶ ማወቅ አልቻሉም ይሆናል።

\end{itemize}

\end{document}
